\subsection{分次模的维数}

在本节中，我们用 M 表示一个分次 H-模（Z 型）。

则 \(\mathbf{S}^{-1}\mathbf{M}\) 是一个 \(\mathbf{S}^{-1}\mathbf{H}\)-模；若置 \(\mathbf{M}_{\geqslant n} = \bigoplus_{i\geqslant n}\mathbf{M}_i\) ，则如

第 1 号中所述，集合序列 \(S^{-1}M_{\geqslant n}\) 是定义在

\(S^{-1}M\) 上的一个穷尽且分离的滤过；且典范映射 \(M \to S^{-1}M\) 诱导出从 M 到分次模 \(\oplus S^{-1}M_{\geqslant n}/S^{-1}M_{\geqslant n+1}\) 的一个同构。

引理 3. --- 假设 H 由 \(\mathbf{H}_{0} \cup \mathbf{H}_{1}\) 生成，而 M 由 \(\bigoplus_{i \leqslant n_{0}} \mathbf{M}_{i}\) 生成，其中 \(n_{0}\) 是适当的整数。则 \((\mathbf{S}^{-1}\mathbf{M}_{\geqslant n})\) 作为 \(\mathbf{S}^{-1}\mathbf{M}\) 的滤过，对于 \(\mathbf{S}^{-1}\mathbf{H}_{\geqslant 1}\) 的理想 \(\mathbf{S}^{-1}\mathbf{H}\) 是良好的。

当 \(n \geqslant n_0\) 时，\(\mathbf{M}_{\geqslant n+1} = \mathbf{H}_1 \cdot \mathbf{M}_{\geqslant n}\)，因而 \(\mathbf{S}^{-1} \mathbf{M}_{\geqslant n+1} = \mathbf{H}_1 \cdot \mathbf{S}^{-1} \mathbf{M}_{\geqslant n} = \mathbf{S}^{-1} \mathbf{H}_{\geqslant 1} \cdot \mathbf{S}^{-1} \mathbf{M}_{\geqslant n}\)。

命题 4. --- 假设 H 是诺特环且 M 是有限生成的。则 \(\dim_{\mathbf{H}}(\mathbf{M}) = \dim_{\mathbf{S}^{-1}\mathbf{H}}(\mathbf{S}^{-1}\mathbf{M})\)。

令 a 为 H-模 M 的湮灭子；它是 H 的一个分次理想。由于 M 是有限生成的 H-模，S \(^{-1}\) H-模 S \(^{-1}\) M 的湮灭子是 S \(^{-1}\) H 的理想 S \(^{-1}\) a。我们有 dim \(_{H}\) (M) = dim(H/a)，且 dim \(_{S^{-1}H}\) (S \(^{-1}\) M) = dim(S \(^{-1}\) H/S \(^{-1}\) a)。命题 4 由第 2 号的定理 1 c) 应用于分次环 H/a 即得。

命题 5. --- 假设 \(\mathbf{H}_{0}\) 是局部阿廷环，\(\mathbf{H}\) 由 \(\mathbf{H}_{0} \cup \mathbf{H}_{1}\) 生成，\(\mathbf{H}_{1}\) 是有限型的 \(\mathbf{H}_{0}\)-模，而 \(\mathbf{M}\) 是非零的有限型 \(\mathbf{H}\)-模。则对于每个 \(\mathbf{M}_{n}\)，\(\mathbf{H}_{0}\) 都是有限生成的 \(n\)-模，具有有限长度且非零；并且存在 \(\mathbf{Q}(\mathbf{T}) \in \mathbf{Z}[\mathbf{T}, \mathbf{T}^{-1}]\)，使得 \(\mathbf{Q}(1) > 0\)，且在环 \(\mathbf{Z}((\mathbf{T}))\) 中有

\[
\sum_ {n \in \mathbf {Z}} \mathrm{long} _ {\mathrm{H} _ {0}} (\mathrm{M} _ {n}). \mathrm{T} ^ {n} = (1 - \mathrm{T}) ^ {- d}. \mathrm{Q} (\mathrm{T}),
\]

其中 \(d = \dim_{\mathbf{H}}(\mathbf{M})\)。

环 \(\mathbf{S}^{-1}\mathbf{H}\) 是局部诺特环（第 1 号，注 1）；\(\mathbf{S}^{-1}\mathbf{H}\)-模 \(\mathbf{S}^{-1}\mathbf{M}\) 是非零有限型模，并且其维数为 \(d = \dim_{\mathbf{H}}(\mathbf{M})\)（命题 4）。此外，\(\mathbf{S}^{-1}\mathbf{H}_{\geqslant 1}\) 是 \(\mathbf{S}^{-1}\mathbf{H}\) 的一个定义理想（§ 3，第 2 号，引理 2），而 \(\mathbf{S}^{-1}\mathbf{M}_{\geqslant n}\) 是 \(\mathbf{S}^{-1}\mathbf{H}_{\geqslant 1}\) 上关于 \(\mathbf{S}^{-1}\mathbf{M}\) 的良好滤过（引理 3）。最后，等式 \(_{\mathbf{S}^{-1}\mathbf{H}}\mathbf{S}^{-1}\mathbf{M}_{\geqslant n}/\mathbf{S}^{-1}\mathbf{M}_{\geqslant n+1} = \text{long}_{\mathbf{H}_0}(\mathbf{M}_n)\) 对每个 \(n\) 成立。因此只需应用第 § 4 节（第 3 号和第 4 号）的定理 2 和定理 3。

注。--- 除了确定整数 d 之外，命题 5 直接由第 § 4 节第 2 号的定理 1 得出。

推论。--- 设 A 是诺特局部环，q 是 A 的一个定义理想。则有 \(\dim(A) = \dim(\text{gr}_{\mathfrak{q}}(A))\)。

将命题 5 应用于 M = H = gr\_q(A) 的情形，得到关系

\[
\sum_ {n \geqslant 0} \mathrm{long} _ {\mathrm{A} / \mathfrak {q}} (\mathfrak {q} ^ {n} / \mathfrak {q} ^ {n + 1}). \mathrm{T} ^ {n} = (1 - \mathrm{T}) ^ {- d} \mathrm{Q} (\mathrm{T})
\]

其中 \(d = \dim(\mathrm{gr}_{\mathfrak{q}}(\mathbf{A}))\) 且 \(\mathbf{Q}(1) \neq 0\)。根据第 § 4 节第 4 号的定理 3，有 \(d = \dim(\mathbf{A})\)，故推论得证。

命题 6. --- 假设 \(\mathbf{H}_{0}\) 是局部阿廷环，\(\mathbf{H}\) 是有限型的 \(\mathbf{H}_{0}\)-代数，而 \(\mathbf{M}\) 是有限型的 \(\mathbf{H}\)-模。

\begin{enumerate}
\def\labelenumi{\alph{enumi})}
\item
  设 \(a_1, \ldots, a_n\) 是 H 中的元素，均为次数 \(> 0\) 的齐次元素；令 \(\varphi\) 为一个（\(\mathbf{H}_0\)-代数）同态，它从 \(\mathbf{H}_0[\mathbf{X}_1, \ldots, \mathbf{X}_n]\) 到 H，并将 \(\mathbf{X}_i\) 送到 \(a_i\)，其中 \(1 \leqslant i \leqslant n\)。则 \(S^{-1}H\)-模 \(S^{-1}M / \sum_{i=1}^{n}(a_i/1) \cdot S^{-1}M\) 具有有限长度，当且仅当 \(\varphi_*(M)\) 是 \(\mathbf{H}_0[\mathbf{X}_1, \ldots, \mathbf{X}_n]\) 上的有限生成模。
\item
  存在一个 H 中元素的族 \((a_{1},...,a_{d})\)，它们都是同一次数 \textgreater0 的齐次元素，其中 \(d=\dim_{\mathrm{H}}(\mathbf{M})\)，并且 \((a_{1}/1,...,a_{d}/1)\) 是 \(S^{-1}H\)-模 \(S^{-1}M\) 的一个极大正割序列。若进一步，H 由 \(H_{1}\) 生成（作为 \(H_{0}\)-代数），且 \(H_{0}\) 的剩余域是无限的，则可取次数为 1 的 \(a_{i}\)。
\item
  置 \(\mathbf{N} = \mathbf{M} / \sum_{i=1}^{n} a_i \mathbf{M}\)。根据命题 4，有 \(\dim_{\mathrm{H}}(\mathbf{N}) = \dim_{\mathrm{S}^{-1}\mathrm{H}}(\mathrm{S}^{-1}\mathrm{N})\)。因此，\(\mathrm{S}^{-1}\mathrm{H}\)-模 \(\mathrm{S}^{-1}\mathrm{N}\) 具有有限长度，当且仅当 H-模 N 具有有限长度；也就是说，当且仅当 N 是有限型的 \(\mathrm{H}_0\)-模。若 \(\varphi_*(\mathbf{M})\) 是定义在 \(\mathrm{H}_0[\mathrm{X}_1, ..., \mathrm{X}_n]\) 上的、通过同态 \(\varphi : \mathrm{H}_0[\mathrm{X}_1, ..., \mathrm{X}_n] \to \mathbf{H}\) 从 M 导出的模，则有 \(\mathbf{N} = \varphi_*(\mathbf{M}) / \sum_{i=1}^{n} X_i \cdot \varphi_*(\mathbf{M})\)。因此（A, II, p.~171，推论 3 和注），\(\varphi_*(\mathbf{M})\) 是 \(\mathrm{H}_0[\mathrm{X}_1, ..., \mathrm{X}_n]\) 上的有限生成模，当且仅当 N 是有限型的 \(\mathrm{H}_0\)-模。这就证明了 a)。
\end{enumerate}

为证明 b)，我们先建立一个引理。

引理 4. --- 设 \(b\) 是 H 中一个次数 \(>0\) 的齐次元素，并且不属于任何极小元 \(\mathfrak{p}\)，这里的极小元取自 \(\operatorname{Supp}(\mathbf{M})\) 且满足 \(\dim(\mathrm{H}/\mathfrak{p}) = \dim_{\mathrm{H}}(\mathbf{M})\)。则 \(\dim_{\mathrm{H}}(\mathbf{M}/b\mathbf{M}) = \dim_{\mathrm{H}}(\mathbf{M}) - 1\)。

置 \(d = \dim_{\mathbf{H}}(\mathbf{M})\)。根据 \(b\) 的定义，有 \(\dim_{\mathbf{H}}(\mathbf{M}/b\mathbf{M}) < d\)。根据命题 4，有

\[
\dim_ {\mathrm{H}} (\mathrm{M} / b \mathrm{M}) = \dim_ {\mathrm{S} ^ {- 1} \mathrm{H}} \left(\mathrm{S} ^ {- 1} \mathrm{M} / (b / 1). \mathrm{S} ^ {- 1} \mathrm{M}\right)
\]

最后，有

\[
\dim_ {\mathrm{S} ^ {- 1} \mathrm{H}} (\mathrm{S} ^ {- 1} \mathrm{M} / (b / 1). \mathrm{S} ^ {- 1} \mathrm{M}) \geqslant \dim_ {\mathrm{S} ^ {- 1} \mathrm{H}} (\mathrm{S} ^ {- 1} \mathrm{M}) - 1.
\]

最后，有

\[
\dim_ {S ^ {- 1} H} (S ^ {- 1} M) = \dim_ {H} (M) = d
\]

因此 \(\dim_{\mathrm{H}}(\mathrm{M}/b\mathrm{M}) \geqslant d - 1\)，从而引理 4 得证。

我们继续证明命题 6 b)。可设 \(\dim_{\mathbf{H}}(\mathbf{M}) > 0\)。注意，\(\operatorname{Supp}(\mathbf{M})\) 的每个极小元都是分次的（将第 2 号的引理 1 应用于 H 除以 M 的湮灭子所得的商）。根据 III, § 1, 第 4 号的命题 8，存在 H 的一个齐次元素 \(b\)，次数为 \(>0\)，且不属于任何极小元 \(\mathfrak{p}\)；这些极小元取自 \(\operatorname{Supp}(\mathbf{M})\)，并满足 \(\dim(\mathrm{H}/\mathfrak{p}) = \dim_{\mathbf{H}}(\mathbf{M})\)。根据引理 4，有 \(\dim_{\mathbf{H}}(\mathbf{M}/b\mathbf{M}) = \dim_{\mathbf{H}}(\mathbf{M}) - 1\)。再假设 H 由 \(\mathrm{H}_{1}\) 生成（作为 \(\mathrm{H}_{0}\)-代数），且剩余域 \(k\) 是 \(\mathrm{H}_{0}\) 的无限域。对于每个满足以下条件的极小元 \(\mathfrak{p}\)：它属于 \(\operatorname{Supp}(\mathbf{M})\)，且满足 \(\dim(\mathrm{H}/\mathfrak{p}) = \dim_{\mathbf{H}}(\mathbf{M})\)，考虑向量子空间 \(\mathrm{V}_{\mathfrak{p}} = (\mathfrak{p} \cap \mathrm{H}_{1}) \otimes_{\mathrm{H}_{0}} k\)，它是 \(k\)-向量空间 \(\mathrm{V} = \mathrm{H}_{1} \otimes_{\mathrm{H}_{0}} k\) 的子空间。若有 \(\mathrm{V}_{\mathfrak{p}} = \mathrm{V}\)，则有 \(\mathfrak{p} \cap \mathrm{H}_1 = \mathrm{H}_1\)（II, §3, 第 2 号，命题 4），从而 \(\mathrm{H}_1 \subset \mathfrak{p}\)，且 \(\dim_{\mathrm{H}}(\mathbf{M}) = \dim(\mathrm{H}/\mathfrak{p}) \leqslant \dim(\mathrm{H}/\mathrm{H}_{\geqslant 1}) = 0\)，这与假设不符。由于假设 \(k\) 是无限的，各个 \(\mathrm{V}_{\mathfrak{p}}\) 的并集不同于 V；若 \(b \in \mathrm{H}_1\) 且 \(b \otimes 1\) 不属于其中任何一个 \(\mathrm{V}_{\mathfrak{p}}\)，则 \(\dim(\mathrm{M}/b\mathrm{M}) = \dim_{\mathrm{H}}(\mathrm{M}) - 1\)。

对 \(d = \dim_{\mathbf{H}}(\mathbf{M})\) 作归纳，便可构造 H 中元素的序列 \((b_1, \ldots, b_d)\)，其中 \(b_i\) 是次数 \(n_i > 0\) 的齐次元素，并且 \(\mathbf{M} / \sum_{i=1}^{n} b_i \mathbf{M}\) 是有限长度的 H-模。若假设 H 由 \(\mathbf{H}_1\) 生成（作为 \(\mathbf{H}_0\)-代数），且 \(\mathbf{H}_0\) 的剩余域是无限的，则可设 \(n_i = 1\)，其中 \(i = 1, \ldots, d\)。根据命题 4，有 \(\dim_{\mathbf{S}^{-1}\mathbf{H}}(\mathbf{S}^{-1}\mathbf{M}) = d\)，并且

\[
\dim_ {\mathrm{S} ^ {- 1} \mathrm{H}} (\mathrm{S} ^ {- 1} \mathrm{M} / \sum_ {i = 1} ^ {d} (b _ {i} / 1). \mathrm{S} ^ {- 1} \mathrm{M}) = 0.
\]

于是 \((b_{1} / 1,\dots,b_{d} / 1)\) 是 \(\mathbf{S}^{-1}\mathbf{H}\)-模 \(\mathbf{S}^{-1}\mathbf{M}\) 的一个极大正割序列。置 \(a_{i} = b_{i}^{(n_{1}\dots n_{d}) / n_{i}}\)，其中 \(1\leqslant i\leqslant d\)；则 \(a_{i}\) 具有相同的次数，并且 \((a_{1} / 1,\dots,a_{d} / 1)\) 是 \(\mathbf{S}^{-1}\mathbf{M}\) 的一个极大正割序列（§ 3，第 2 号，注 3）。

推论 1. --- 假设 \(\mathbf{H}_0\) 是一个域，且 \(\mathbf{H}\) 是有限型的 \(\mathbf{H}_0\)-代数。置 \(n = \dim(\mathbf{H})\)。存在齐次元素 \(a_1, ..., a_n\)，它们属于 \(\mathbf{H}\)，且都具有 \(>0\) 的同一次数，使 \(\mathbf{H}_0\)-同态 \(\varphi: \mathbf{H}_0[\mathbf{X}_1, ..., \mathbf{X}_n] \to \mathbf{H}\)（由 \(\varphi(\mathbf{X}_i) = a_i\)，\(i = 1, ..., n\) 定义）是单射，并使 \(\mathbf{H}\) 成为有限的 \(\mathbf{H}_0[\mathbf{X}_1, ..., \mathbf{X}_n]\)-代数。若 \(\mathbf{H}\) 由 \(\mathbf{H}_1\) 生成（作为 \(\mathbf{H}_0\)-代数），且 \(\mathbf{H}_0\) 是无限的，则可取次数为 1 的 \(a_i\)。

根据命题 6，存在一个上述形式的 \(\mathbf{H}_0\)-同态 \(\varphi\)，使 H 成为 \(\mathbf{H}_0[\mathbf{X}_1, \ldots, \mathbf{X}_n]\)-有限代数。根据第 § 2 节第 3 号的定理 1，有

\[
\dim \left(\mathrm{H} _ {0} \left[ \mathrm{X} _ {1}, \dots , \mathrm{X} _ {n} \right] / (\operatorname{Ker} \varphi)\right) = \dim (\mathrm{H});
\]

因为有

\[
\dim (\mathrm{H}) = n = \dim \left(\mathrm{H} _ {0} \left[ \mathrm{X} _ {1}, \dots , \mathrm{X} _ {n} \right]\right),
\]

且 \(\mathrm{H}_0[\mathrm{X}_1, \dots, \mathrm{X}_n]\) 是整环，这就推出 \(\operatorname{Ker} \varphi = \{0\}\)。

注。--- 设 \((h_{1}, \ldots, h_{r})\) 是 \(H_{0}\)-向量空间 \(H_{1}\) 的一个有限生成系。对于 \(\lambda = (\lambda_{1}, \ldots, \lambda_{r}) \in H_{0}^{r}\)，置 \(h_{\lambda} = \lambda_{1} h_{1} + \cdots + \lambda_{r} h_{r}\)。命题 6 和推论 1 的证明给出如下结果：由 \((\lambda_{1}, \ldots, \lambda_{n})\) 组成的、属于 \((\mathrm{H}_{0}^{r})^{n}\) 且使元素 \(a_{i} = h_{\lambda_{i}} \in H_{1}\) 满足推论 1 结论的元素集合，包含 \((\mathrm{H}_{0}^{r})^{n}\) 中有限个不等于整个空间的向量子空间之并的补集。

推论 2. --- 设 A 是诺特局部环且 \(n\in\mathbb{N}\)。要有 \(\dim(\mathbf{A})\geqslant n\)，充要条件是对于每个整数 \(r\geqslant0\)，有

\[
\left[ \mathfrak {m} _ {\mathrm{A}} ^ {r} / \mathfrak {m} _ {\mathrm{A}} ^ {r + 1}: \kappa_ {\mathrm{A}} \right] \geqslant \binom {n + r - 1} {n - 1}, \quad \left(\text { resp.   long } _ {\mathrm{A}} (\mathrm{A} / \mathfrak {m} _ {\mathrm{A}} ^ {r + 1}) \geqslant \binom {n + r} {n}\right);
\]

对于每个 r 都取等号，当且仅当 A 是维数为 n 的正则环。

该条件是充分的（§ 4，第 4 号，定理 3 以及 § 5，第 2 号，定理 1）。下面证明其必要性。考虑分次环 gr(A) = gr \(_{m_A}\) (A)；令 k 为 κ \(_A\) 的一个无限扩域，并置 H = k ⊗ κ \(_A\) gr(A)。环 H 的维数 ≥ n（命题 5 及其推论）；因此由推论 1，得到分次 k-代数的一个单射分次同态 φ : H \(_0\) {[}X \(_1\) , \ldots, X \(_n\) {]} → H。于是，对于每个整数 r ≥ 0，有

\[
\left[ \mathrm{gr} _ {r} (\mathrm{A}): \kappa_ {\mathrm{A}} \right] = \left[ \mathrm{H} _ {r}: \mathrm{H} _ {0} \right] \geqslant \binom {n + r - 1} {n - 1},
\]

且每个 \(r\) 都取等号会推出 \(\varphi\) 为双射，从而 A 正则（ \(\S\) 5，第 2 号，定理 1）。

等式

\[
\mathrm{long} _ {\mathbf {A}} (\mathrm{A} / \mathfrak {m} _ {\mathbf {A}} ^ {r + 1}) = \sum_ {i = 0} ^ {r} \left[ \mathrm{gr} _ {i} (\mathbf {A}): \kappa_ {\mathbf {A}} \right]
\]

以及

\[
\binom {n + r} {n} = \sum_ {i = 0} ^ {r} \binom {n + i - 1} {n - 1}
\]

从而得到关于函数 \(r \mapsto \text{long}_{\mathbf{A}}(\mathbf{A}/m_{\mathbf{A}}^{r+1})\) 的相应断言。

